\DocumentMetadata{
  pdfversion=2.0,
  pdfstandard=ua-2,
  lang=en-US,
  tagging=on,
  tagging-setup={math/setup=mathml-SE,tabsorder=structure}
}
\documentclass[11pt,letterpaper]{article}
\usepackage[margin=0.68in,includefoot]{geometry}
\usepackage{newcomputermodern}
\usepackage{amsmath,amscd,mathtools}
\usepackage{tikz,adjustbox}
\providecommand{\square}{\mdlgwhtsquare}
\usepackage[hidelinks]{hyperref}
\hypersetup{
  pdftitle={Logic PhD Exam, May 2010},
  pdfauthor={University of Florida Department of Mathematics},
  pdfsubject={Graduate mathematics examination},
  pdfdisplaydoctitle=true,
  bookmarksopen=true
}
\setcounter{secnumdepth}{0}
\setlength{\parindent}{0pt}
\setlength{\parskip}{0.28em}
\setlength{\emergencystretch}{3em}
\setlength{\leftmargini}{2.15em}
\setlength{\leftmarginii}{2.65em}
\setlength{\leftmarginiii}{3em}
\pagestyle{empty}
\raggedbottom
\allowdisplaybreaks
\NewTaggingSocketPlug{block/list/label}{examproblem}{%
  \tagstructbegin{tag=Lbl}\tagstructend%
  \tagstructbegin{tag=\LBody}%
  \tagstructbegin{tag=\ProblemHeadingTag,title-o={\ProblemHeadingText}}%
  \tagmcbegin{tag=\ProblemHeadingTag,actualtext={\ProblemHeadingText}}%
  #2%
  \tagmcend%
  \tagstructend%
}
\newenvironment{examproblems}[2]{%
  \def\ProblemHeadingTag{#2}%
  \AssignTaggingSocketPlug{block/list/label}{examproblem}%
  \begin{enumerate}%
  \setcounter{enumi}{#1}%
  \renewcommand{\labelenumi}{\textbf{\theenumi.}}%
  \setlength{\topsep}{0.35em}%
  \setlength{\partopsep}{0pt}%
  \setlength{\itemsep}{0.48em}%
  \setlength{\parsep}{0.18em}%
}{\end{enumerate}}
\newenvironment{examsubparts}{%
  \AssignTaggingSocketPlug{block/list/label}{default}%
  \begin{enumerate}%
  \setlength{\topsep}{0.18em}%
  \setlength{\partopsep}{0pt}%
  \setlength{\itemsep}{0.16em}%
  \setlength{\parsep}{0.1em}%
}{\end{enumerate}}
\newenvironment{examinstructions}{%
  \par\begingroup\itshape
}{%
  \par\endgroup\vspace{0.18em}
}
\makeatletter
\renewcommand\subsection{\@startsection{subsection}{2}{0pt}%
  {-1.25ex plus -.25ex minus -.1ex}%
  {0.55ex plus .1ex}%
  {\normalfont\large\bfseries}}
\renewcommand{\@maketitle}{%
  \newpage\null\vskip -1em%
  {\footnotesize\bfseries
    \href{https://www.ufl.edu/}{University of Florida}, \href{https://math.ufl.edu/}{Department of Mathematics}\par}%
  \vskip -0.5em%
  \tagstructbegin{tag=H1,title={Logic PhD Exam, May 2010}}%
  \tagpdfparaOff%
  \tagmcbegin{tag=H1}%
    {\LARGE\bfseries\href{https://gma.math.ufl.edu/exams/logic-phd/}{Logic PhD Exam}, May 2010\par}%
  \tagmcend%
  \tagpdfparaOn%
  \tagstructend%
  \vskip 0.25em%
  \tagmcbegin{artifact}%
    \hrule height 1pt%
  \tagmcend%
  \vskip 1em%
}
\makeatother
\title{Logic PhD Exam, May 2010}
\author{}
\date{}
\begin{document}
\maketitle
\thispagestyle{empty}
\begin{examinstructions}
Answer six questions and at least one from each section.
\end{examinstructions}
\subsection{Section 1}
\begin{examproblems}{0}{H3}
\def\ProblemHeadingText{Problem 1.}
\item
Sketch the proof of Gödel’s Completeness Theorem for Predicate Logic and explain the use of Henkin constants (or Skolem functions).
\def\ProblemHeadingText{Problem 2.}
\item
Show that the theory of the class of infinite sets is complete and decidable, but is not finitely axiomatizable.
\def\ProblemHeadingText{Problem 3.}
\item
Sketch a proof of the Vaught test for categoricity using the Löwenheim–Skolem Theorem.
\def\ProblemHeadingText{Problem 4.}
\item
Show that if~\(A\) and~\(B\) are two countable atomic structures which are elementarily equivalent, then \(A \simeq_p B\).
\end{examproblems}
\subsection{Section 2}
\begin{examproblems}{4}{H3}
\def\ProblemHeadingText{Problem 5.}
\item
(ZFC) For any infinite cardinal~\(\kappa\), \(\kappa \times \kappa = \kappa\).
\def\ProblemHeadingText{Problem 6.}
\item
Describe the constructible universe~\(L\) and sketch a proof that~\(L\) satisfies the Continuum Hypothesis.
\def\ProblemHeadingText{Problem 7.}
\item
Show that for any set~\(u\) and any countable set~\(v\), \(FF(u,v)\) satisfies the countable (anti)-chain condition.
\end{examproblems}
\subsection{Section 3}
\begin{examproblems}{7}{H3}
\def\ProblemHeadingText{Problem 8.}
\item
Sketch a proof of Gödel’s incompleteness Theorem for Arithmetic. –You may assume the Representability Hypothesis.
\def\ProblemHeadingText{Problem 9.}
\item
Construct a simple r.e. set.
\def\ProblemHeadingText{Problem 10.}
\item
Show that the following are equivalent for any infinite \(X \subseteq \omega\):
\begin{examsubparts}
\item[\textnormal{(a)}]
\(X\) is recursively enumerable (that is, the range of a total recursive function);
\item[\textnormal{(b)}]
\(X\) is semirecursive (there is a recursive relation~\(R\) such that, for all~\(x\), \(x \in X \Longleftrightarrow (\exists y)R(x,y)\).
\item[\textnormal{(c)}]
\(X\) is the domain of a partial recursive function.
\end{examsubparts}
\end{examproblems}
\end{document}
